\documentclass{training}
\schoollogo{../../../resources/latex/crack.pdf}
\usepackage[utf8]{inputenc}
\usepackage[T1]{fontenc}
\usepackage{lmodern}
\setlang{french}
\usepackage{hyperref}
\usepackage{graphicx}
\usepackage{tabularx}
\usepackage{fvextra}
\fvset{breaklines=true,breakanywhere=true,breaksymbol={},breaksymbolleft={},breaksymbolright={},breakanywheresymbolpre={},breakanywheresymbolpost={}}
\makeatletter
\g@addto@macro\UrlBreaks{%
  \do\a\do\b\do\c\do\d\do\e\do\f\do\g\do\h\do\i\do\j\do\k\do\l\do\m%
  \do\n\do\o\do\p\do\q\do\r\do\s\do\t\do\u\do\v\do\w\do\x\do\y\do\z%
  \do\A\do\B\do\C\do\D\do\E\do\F\do\G\do\H\do\I\do\J\do\K\do\L\do\M%
  \do\N\do\O\do\P\do\Q\do\R\do\S\do\T\do\U\do\V\do\W\do\X\do\Y\do\Z%
  \do\0\do\1\do\2\do\3\do\4\do\5\do\6\do\7\do\8\do\9%
  \do\<\do\>\do\(\do\)\do\'\do\/}
\schoollogo{../../../resources/latex/crack.pdf}
\makeatother


\begin{document}

\title{Windows : réflexion et frontières --- A-EPI-004}
\subtitle{Réflexion, frontières et correctifs incomplets}
\member{Ted ``Tedsig42'' Kouhouenou}{ted.sig42@tutamail.com}
\summary{Module avancé : modéliser la chaîne d'authentification Windows comme une machine à états, puis analyser une famille de CVE récentes qui partagent une racine de réflexion, de contournement de frontière ou de chargement privilégié. Le laboratoire est un AD simulé, jamais une cible réelle.}
\version{1.0}
\prerequisite{A-EPI-001, bases d'Active Directory, NTLM, Kerberos, SMB}
\maketitle

\newpage
\tableofcontents
\newpage

\chapter{Instructions}
\begin{itemize}
  \item Trois jours.
  \item Aucun test sur un domaine de production ou un domaine tiers.
  \item Le laboratoire est un environnement Active Directory simulé,
        isolé, avec des comptes jetables.
  \item Aucune coercition d'authentification, aucun relais réel, aucun
        outil offensif contre une cible non autorisée.
  \item Les CVE de 2026 doivent être traitées comme des cas d'étude :
        vérifiez chaque affirmation dans les avis publics.
  \item Dark Elevator et CVE-2026-66804 sont des cas de logique et de
        correctif incomplet : aucune reproduction d'exploit n'est demandée.
\end{itemize}

\subsection*{Format du rendu --- en-tête obligatoire}

Tout fichier rendu dans ce module doit commencer par l'en-tête
ci-dessous, \textbf{adapté à la syntaxe de commentaire du langage}
utilisé.

\subsubsection*{Modèle générique}

\begin{minted}{c}
/**
 * Module: <code module>
 * Exercise: <numéro exo>
 * File: <nom du fichier>
 * Author: <votre login JeuneCrack>
 *
 * Description: <une phrase, pas de roman>
 *
 * Return: <voir ci-dessous>
 */
\end{minted}

\subsubsection*{Login}

Votre identifiant \textbf{JeuneCrack} (ex: \texttt{tedsig42}).
Pour un rendu en binôme : \texttt{login1/login2}.
Ajoutez un handle Matrix après le login si vous le souhaitez
(ex: \texttt{tedsig42 @tedsig42}).

\subsubsection*{Return --- selon le type de fichier}

\begin{tabularx}{\linewidth}{lX}
  \toprule
  \textbf{Type de fichier} & \textbf{Valeur de Return} \\
  \midrule
  C (\texttt{.c}) &
  \texttt{return (0);} (succès),
  \texttt{return (-1);} (erreur),
  \texttt{return (valeur);} (résultat calculé) \\

  Python (\texttt{.py}) &
  \texttt{0} (succès),
  \texttt{1} (erreur),
  \texttt{N/A} (analyse, script sans exit) \\

  Shell (\texttt{.sh}) &
  \texttt{exit 0} (succès),
  \texttt{exit 1} (erreur) \\

  Markdown (\texttt{.md}) &
  \texttt{N/A} --- les fichiers d'analyse n'ont pas de valeur de retour \\

  Exploit C/Python &
  \texttt{shell} (shell obtenu),
  \texttt{flag: <flag>} (flag capturé),
  \texttt{proof: <hash>} (preuve de concept) \\

  Binaire (\texttt{.bin}) &
  \texttt{N/A} --- pas d'en-tête pour les fichiers binaires \\
  \bottomrule
\end{tabularx}

\subsubsection*{Exemples par langage}

\textbf{C (\texttt{.c})} :
\begin{minted}{c}
/**
 * Module: F-EXP-001
 * Exercise: ex02
 * File: exploit.c
 * Author: tedsig42
 *
 * Description: Buffer overflow with ret2libc
 *
 * Return: 0
 */
\end{minted}

\textbf{Python (\texttt{.py})} :
\begin{minted}{python}
# Module: F-WEB-001
# Exercise: ex05
# File: ssrf-exploit.py
# Author: tedsig42
#
# Description: Blind SSRF bypass via redirect
#
# Return: 0
\end{minted}

\textbf{Shell (\texttt{.sh})} :
\begin{minted}{bash}
# Module: F-SUP-001
# Exercise: ex03
# File: dependency-confusion.sh
# Author: tedsig42
#
# Description: Automated dependency confusion test
#
# Return: exit 0
\end{minted}

\textbf{Markdown (\texttt{.md})} --- analyse, rapport :
\begin{Verbatim}[fontsize=\normalsize]
<!--
Module: M-SRC-001
Exercise: ex04
File: audit-report.md
Author: tedsig42

Description: Code audit of vulnerable-httpd

Return: N/A
-->
\end{Verbatim}

\chapter{Instructions Communes}

\section{Structure du D\'ep\^ot}

Votre d\'ep\^ot Git doit respecter la convention de nommage suivante pour
tous les modules soumis. L'outil d'\'evaluation (\texttt{pendragon})
se base sur cette structure pour localiser et tester votre travail.

\begin{console}
<login>/
+-- <MODULE-ID>-<slug>/         # ex. M-BLS-001-business-logic/
|   +-- <turnindir>/            # ex. ex00-pricing-manipulation/
|   |   +-- <file1>             # ex. poc-pricing.py
|   |   +-- <file2>             # ex. report.md
|   |   +-- ...
|   +-- <turnindir>/            # ex. ex01-race-conditions/
|   |   +-- ...
|   +-- ...
+-- ...
\end{console}

\begin{itemize}
  \item \textbf{R\'epertoire du module} : le nom complet du r\'epertoire tel
        qu'affich\'e dans votre curriculum (ex. \texttt{M-BLS-001-business-logic}).
  \item \textbf{R\'epertoire de rendu} : la valeur de \texttt{\textbackslash turnindir}
        fournie dans chaque exercice (ex. \texttt{ex00-pricing-manipulation}).
  \item \textbf{Fichiers de l'exercice} : exactement les fichiers list\'es dans
        \texttt{\textbackslash exfiles}. Ni fichiers suppl\'ementaires, ni fichiers manquants.
\end{itemize}

\section{Qu'est-ce qu'une Preuve de Concept Valide}

Une PoC doit \^etre :

\begin{enumerate}
  \item \textbf{Reproductible}. N'importe qui avec le m\^eme environnement peut
        l'ex\'ecuter et observer la vuln\'erabilit\'e. Les chemins cod\'es en dur,
        les adresses IP absolues et les jetons expir\'es ne sont pas acceptables.
  \item \textbf{Document\'ee}. Chaque \'etape est expliqu\'ee. L'\'evaluteur ne
        devrait pas avoir \`a deviner ce que fait le script.
  \item \textbf{Cibl\'ee}. L'exercice sp\'ecifie une cible autoris\'ee
        (lab, maquette, binaire local ou plateforme d\'esign\'ee). Ne testez
        jamais les PoC sur des syst\`emes de production qui ne vous appartiennent pas.
  \item \textbf{Minimale}. La PoC d\'emontre \emph{uniquement} la vuln\'erabilit\'e
        en question. Les charges utiles non li\'ees et le bruit sont p\'enalis\'es.
\end{enumerate}

\begin{warning}
  Une PoC qui modifie silencieusement l'\'etat du syst\`eme, supprime des
  fichiers ou persiste des changements en dehors du p\'erim\`etre de l'exercice
  sera rejet\'ee. Si votre exploit n\'ecessite un nettoyage, documentez les
  \'etapes de nettoyage dans votre rapport.
\end{warning}

\section{Processus de Soumission}

\begin{enumerate}
  \item R\'ealisez les livrables de l'exercice sur votre machine locale.
  \item Committez et poussez vers votre d\'ep\^ot Git assign\'e.
  \item Le d\'ep\^ot est gel\'e apr\`es la date limite.
  \item Les pairs clonent votre d\'ep\^ot et \'evaluent en utilisant :
        \begin{itemize}
          \item \texttt{pendragon} pour l'ex\'ecution automatis\'ee des tests
                (l\`a o\`u le module fournit des configurations mouli).
          \item La grille manuelle pour la qualit\'e du rapport, le style de
                code et la m\'ethodologie.
        \end{itemize}
\end{enumerate}

\hint{Ex\'ecutez toujours \texttt{pendragon run <MODULE-ID> <ex\_id>
<path/to/student/repo>} avant votre commit final pour v\'erifier
que votre soumission passe les v\'erifications automatis\'ees.}

\subsection*{Test avec pendragon}

Si le module inclut une configuration mouli (fichiers JSON dans
\texttt{resources/mouli/}), vous pouvez tester votre soumission localement :

\begin{console}
pendragon run M-BLS-001 ex00 ~/student-repos/tedsig42
\end{console}

\texttt{pendragon} r\'esout le chemin de l'exercice automatiquement en
cherchant \texttt{<mod-fullname>/<ex-dirname>/} dans le r\'epertoire
\'etudiant donn\'e.

\subsection*{V\'erification de style avec caliban}

Pour les modules qui n\'ecessitent du code C, \texttt{caliban} v\'erifie
la conformit\'e du style :

\begin{console}
caliban check ~/student-repos/tedsig42/M-BLS-001-business-logic
\end{console}

\texttt{caliban} applique la Norme 42 (taille des fonctions, nommage
des variables, constructions interdites) adapt\'ee aux contextes de
recherche de vuln\'erabilit\'es. Ex\'ecutez-le avant la soumission ---
un \'echec aux v\'erifications de style peut invalider l'exercice.

\section{R\`egle des Fichiers Suppl\'ementaires Interdits}

\begin{itemize}
  \item Votre d\'ep\^ot doit contenir \textbf{uniquement} les fichiers
        list\'es dans \texttt{\textbackslash exfiles} pour chaque exercice.
  \item Fichiers interdits : \texttt{.DS\_Store}, \texttt{\_\_pycache\_\_},
        \texttt{.venv/}, \texttt{node\_modules/}, \texttt{.git/},
        binaires compil\'es (sauf explicitement list\'es dans \texttt{exfiles}),
        fichiers d'\'echange d'\'editeur, fichiers journaux et tout fichier
        non demand\'e.
  \item Utilisez \texttt{.gitignore} pour exclure automatiquement les
        fichiers ind\'esirables.
\end{itemize}

\section{Format du Rapport}

Lorsqu'un rapport (\texttt{report.md}) est requis, il doit inclure :

\begin{enumerate}
  \item \textbf{Classe de vuln\'erabilit\'e} (identifiant CWE si applicable).
  \item \textbf{Étapes de reproduction} --- num\'erot\'ees, pr\^etes
        \`a \^etre copi\'ees-coll\'ees.
  \item \textbf{Charges utiles} --- chaînes exactes, requ\^etes ou
        entr\'ees utilis\'ees.
  \item \textbf{Impact} --- ce qu'un attaquant peut accomplir.
  \item \textbf{Mitigation} --- comment le d\'eveloppeur devrait corriger.
  \item \textbf{Sources} --- toutes les ressources consult\'ees (voir
        Section~\ref{sec:cite-sources}).
\end{enumerate}

\subsection*{En-t\^ete de Rendu}

Chaque fichier source doit inclure l'en-t\^ete de rendu d\'efini dans
\path{../../../resources/latex/deliverable-header.tex}, adapt\'e
\`a la syntaxe de commentaire du langage utilis\'e.

\section{Cibles Autoris\'ees Uniquement}

\begin{warning}
  Vous pouvez \textbf{uniquement} tester les cibles explicitement list\'ees
  dans l'exercice (labs, serveurs locaux, maquettes, binaires d\'esign\'es).
  Tester des syst\`emes de production r\'eels, des applications tierces
  ou toute cible qui ne vous appartient pas ou pour laquelle vous n'avez
  pas d'autorisation \'ecrite est interdit et peut entraîner une
  disqualification imm\'ediate.
\end{warning}

L'\'ethique du bug bounty s'applique tout au long de ce curriculum :
\begin{itemize}
  \item Aucun test sur des programmes en production sans un accord de
        p\'erim\`etre valide.
  \item Aucun d\'eni de service sauf si l'exercice l'exige explicitement.
  \item Aucune exfiltration de donn\'ees au-del\`a de ce qui est
        n\'ecessaire pour d\'emontrer la vuln\'erabilit\'e.
  \item Si vous acc\'edez accidentellement \`a des donn\'ees r\'eelles
        d'utilisateurs, arr\^etez imm\'ediatement et documentez ce qui
        s'est pass\'e.
\end{itemize}

\section{Citez Toutes Vos Sources}
\label{sec:cite-sources}

Toute ressource consult\'ee pendant votre travail doit \^etre cit\'ee.
Cela inclut :

\begin{itemize}
  \item Articles acad\'emiques et billets de blog.
  \item Rapports publics de bug bounty (HackerOne, Bugcrowd, etc.).
  \item Documentation officielle (RFC, pages de manuel, documentation
        fournisseur).
  \item \textbf{Outils d'IA} : nom du mod\`ele, r\'esum\'e de la
        question, date et m\'ethode de v\'erification (voir le chapitre
        Instructions IA pour le format de citation exact).
  \item Extraits de code adapt\'es de projets open-source (avec
        attribution de licence si requis).
\end{itemize}

\subsection*{Format de citation}

\begin{minted}{markdown}
# R\'ef\'erences

- [1] PortSwigger, "Business Logic Flaws",
  https://portswigger.net/web-security/logic-flaws
  (consult\'e le 2025-03-15).
- [2] Rapport HackerOne #123456, "IDOR on /api/users",
  anonymis\'e, consult\'e via divulgation publique.
- [3] G\'en\'er\'e avec l'assistance de Claude 3.5 Sonnet
  le 2025-03-14. Question : "Write a Python script that...".
  V\'erifi\'e par : revue manuelle de chaque requ\^ete HTTP.
\end{minted}

\subsection*{Attribution en ligne}

Si un exercice vous demande explicitement d'utiliser une technique ou
un outil sp\'ecifique, vous n'avez pas besoin de citer l'exercice
lui-m\^eme comme source. Citez tout ce qui se trouve en dehors de
l'\'enonc\'e de l'exercice.

\chapter{Instructions sur l'IA}

\section{Aperçu de la Politique}

Les outils de codage assistés par IA (ChatGPT, Claude, Copilot, Gemini,
Cursor, Codeium, et similaires) sont \textbf{autorisés} tout au long
de ce cursus. La recherche en sécurité est un domaine où la maîtrise
de l'IA devient essentielle --- apprendre à utiliser ces outils
efficacement fait partie de votre formation.

Cependant, l'objectif est la \textbf{compréhension}, pas l'automatisation.
Si vous ne pouvez pas expliquer chaque ligne de votre exploit sans
l'aide de l'IA, vous n'avez pas terminé l'exercice.

\begin{warning}
  Utiliser l'IA pour contourner l'objectif pédagogique d'un exercice
  --- par exemple, lui demander de rétro-ingéniérer un binaire pour
  ne pas avoir à apprendre le désassembleur, ou de générer une chaîne
  d'exploitation complète que vous ne comprenez pas --- va à l'encontre
  du but du cursus et sera traité comme une fraude académique.
\end{warning}

\section{Utilisations Autorisées}

\begin{itemize}
  \item \textbf{Explication de concepts}. Demander à l'IA de clarifier
        une classe de vulnérabilité, un sous-système du noyau, ou un
        protocole réseau.
  \item \textbf{Échafaudage de code}. Générer du code de départ pour
        un harnais, un fuzzer, ou un squelette de PoC --- puis le
        modifier et l'étendre vous-même.
  \item \textbf{Débogage}. Expliquer des messages d'erreur, la sortie
        d'un désassembleur, ou un comportement de crash inattendu.
  \item \textbf{Résumé de documentation}. Condenser une longue RFC ou
        un article technique en notes exploitables.
  \item \textbf{Traduction}. Traduire du contenu technique entre
        l'anglais et le français (les deux langues sont utilisées
        dans ce cursus).
  \item \textbf{Rédaction de rapports}. Demander de l'aide pour
        structurer un rapport de bug bounty --- mais remplir le
        contenu technique vous-même.
\end{itemize}

\section{Utilisations Interdites}

\begin{itemize}
  \item \textbf{Génération complète d'exploit}. Demander à l'IA
        d'écrire l'exploit complet d'un exercice, puis le soumettre
        sans le comprendre.
  \item \textbf{Contournement des contraintes de l'exercice}. Utiliser
        l'IA pour contourner une restriction que l'exercice impose
        explicitement (p.\ ex.\ «~n'utilisez pas de scanners
        automatisés~»).
  \item \textbf{Découverte automatisée de vulnérabilités}. Exécuter des
        outils de scan pilotés par IA qui remplacent l'analyse manuelle
        que vous êtes censé pratiquer.
  \item \textbf{Rédaction des grilles d'évaluation}. Vous n'êtes pas
        autorisé à utiliser l'IA pour générer des critères d'évaluation
        par les pairs ou des réponses.
  \item \textbf{Plagiat sans attribution}. Copier le résultat de l'IA
        sans le citer (voir Section~\ref{sec:ai-citation}).
\end{itemize}

\section{Responsabilité Personnelle}

Lorsque vous demandez à l'IA de générer du code d'exploit, vous
devenez responsable de chaque octet de ce résultat avant de
l'utiliser contre une cible.

\begin{enumerate}
  \item \textbf{Lisez le code généré} ligne par ligne.
  \item \textbf{Identifiez toutes les hypothèses} faites par l'IA
        (architecture, version du noyau, drapeaux du compilateur,
        topologie réseau).
  \item \textbf{Testez en isolation} avant d'exécuter contre une
        cible d'exercice.
  \item \textbf{Comprenez les modes d'échec}. Si l'exploit généré
        plante ou se bloque, vous devez diagnostiquer et corriger
        le problème vous-même.
\end{enumerate}

\hint{Si l'IA génère du code que vous ne comprenez pas, demandez-lui
de l'expliquer. Si vous ne comprenez toujours pas après
l'explication, le code n'est pas prêt à être utilisé. Passez à
une approche plus simple que vous maîtrisez entièrement.}

\section{Citation Obligatoire de l'IA}
\label{sec:ai-citation}

Chaque fichier contenant du code généré ou assisté par l'IA doit
inclure un commentaire de citation. Le format est~:

\begin{minted}{python}
# Generated with assistance from <tool> on <date>.
# Prompt: <one-line summary of the prompt>
# Verified by: <method used to verify the output>
\end{minted}

\subsection*{Exemples}

\textbf{Python}~:
\begin{minted}{python}
# Generated with assistance from Claude 3.5 Sonnet on 2025-03-14.
# Prompt: "Write a Python script that sends HTTP requests with
# modified price fields to /api/checkout/pay"
# Verified by: manual review of each request and response
\end{minted}

\textbf{C}~:
\begin{minted}{c}
/* Generated with assistance from GitHub Copilot on 2025-03-14.
 * Prompt: "Implement a stack-based buffer overflow with
 * ret2libc for x86_64 Linux"
 * Verified by: compiled with -fno-stack-protector and
 * tested against local binary with ASLR disabled */
\end{minted}

\textbf{Markdown (rapport)}~:
\begin{minted}{html}
<!--
Generated with assistance from ChatGPT-4 on 2025-03-14.
Prompt: "Structure a bug bounty report for IDOR with
sections: summary, steps, impact, mitigation"
Verified by: all technical content written manually,
AI only suggested structure
-->
\end{minted}

\subsection*{Interactions multiples avec l'IA}

Si vous avez utilisé plusieurs outils d'IA, ou itéré à travers
plusieurs invites, citez chaque interaction significative. Un
groupe d'invites connexes peut être résumé~:

\begin{minted}{html}
<!--
AI assistance used in this file:
- Claude 3.5 Sonnet (2025-03-14): harness skeleton generation
- ChatGPT-4 (2025-03-14): crash triage help
- Copilot (2025-03-15): inline autocomplete for loop logic
All generated code reviewed and modified manually.
-->
\end{minted}

\section{Évaluation par les Pairs et l'IA}

Lors de l'évaluation par les pairs, l'évaluateur peut~:

\begin{itemize}
  \item Vous demander d'\textbf{expliquer tout code généré par l'IA}
        ligne par ligne.
  \item Vous demander de \textbf{reproduire le raisonnement} sans
        l'aide de l'IA.
  \item Demander à voir votre \textbf{historique d'invites} pour un
        exercice donné.
  \item Déduire des points si la citation de l'IA est manquante ou
        si vous ne pouvez pas expliquer le code que l'IA a généré.
\end{itemize}

\begin{warning}
  Présenter le code généré par l'IA comme votre propre travail
  (défaut de citation) est du plagiat. Si l'évaluateur soupçonne
  du plagiat, il peut saisir l'équipe pédagogique pour examen.
\end{warning}

\section{Rappel Éthique}

\begin{itemize}
  \item Les modèles d'IA peuvent produire du code incorrect,
        dangereux, ou illégal. Vous êtes \textbf{responsable} de
        vérifier chaque résultat avant exécution.
  \item Les modèles d'IA entraînés sur la recherche publique en
        sécurité peuvent reproduire textuellement du code provenant
        de projets sous licence GPL. Vérifiez les licences avant
        d'intégrer ces éléments dans votre travail.
  \item Ne saisissez pas de code sensible ou propriétaire dans des
        outils d'IA hébergés dans le cloud. Utilisez des modèles
        locaux (Llama, Mistral, DeepSeek) pour les travaux impliquant
        des cibles non publiques.
  \item L'éthique de la recherche en sécurité que vous apprenez ici
        s'applique à l'utilisation de l'IA~: ne demandez pas à l'IA
        de générer des malwares, des rançongiciels, ou du code destiné
        à nuire à des systèmes en dehors du cadre de l'exercice.
\end{itemize}

\info{La maîtrise des outils d'IA est une compétence recherchée dans
la chasse aux bogues moderne. Utilisez-les sagement, citez-les
honnêtement, et conservez toujours votre propre compréhension de
chaque exploit que vous déployez.}


\chapter{Avant-propos}

La chaîne d'authentification Windows n'est pas un seul composant. C'est une
machine à états répartie entre le client, le serveur, le contrôleur de
domaine et les services qui acceptent des tickets ou des jetons.

Une \emph{réflexion} apparaît quand une authentification destinée à une
frontière de confiance est réutilisée sur une autre frontière. Un
\emph{contournement de frontière} apparaît quand un composant accepte une
preuve qui n'a pas été produite pour lui.

\begin{center}
\begin{tabularx}{\linewidth}{|p{0.22\linewidth}|X|}
\hline
\textbf{Étape} & \textbf{Question d'audit} \\ \hline
Preuve & Qui a produit cette preuve, pour quelle cible, et pour quelle
durée ? \\ \hline
Transport & La preuve est-elle liée au canal ou au service qui la reçoit ? \\ \hline
Validation & Le vérificateur applique-t-il la même règle que l'émetteur ? \\ \hline
Translation & La preuve est-elle transformée en un autre type de jeton
sans perte de contexte ? \\ \hline
\end{tabularx}
\end{center}

\begin{warn}
Les CVE de 2026 évoluent. Ne reprenez pas un chiffre ou un mécanisme sans
citer l'avis officiel. Le but du module est la méthode, pas la mémorisation
d'un exploit.
\end{warn}

\chapter{Aperçu du module}

\begin{center}
\begin{tabularx}{\linewidth}{|c|l|c|c|}
\hline
\textbf{Ex} & \textbf{Titre} & \textbf{Difficulté} & \textbf{Points} \\ \hline
00 & Carte de la machine à états & Moyen & 4 \\
01 & CVE-2025-33073 : réflexion NTLM & Difficile & 5 \\
02 & CVE-2025-58726 : Ghost SPNs & Difficile & 4 \\
03 & 2026 : chaîne et correctif incomplet & Difficile & 5 \\
04 & Laboratoire et détection & Moyen & 4 \\ \hline
& \textbf{Total} & & \textbf{22} \\ \hline
\end{tabularx}
\end{center}

\startexercices

\chapter{Exercise 00}
\extitle{Carte de la machine à états}
\turnindir{ex00}
\exnumber{00}
\exfiles{fsm.md}
\exauthorize{un éditeur de texte, documentation Microsoft}
\makeheaderfiles

Dans \texttt{fsm.md}, modélisez une authentification NTLM/Kerberos
simplifiée :

\begin{enumerate}
  \item les acteurs : client, serveur, KDC, service ;
  \item les messages échangés ;
  \item les preuves produites et leur cible ;
  \item les transitions d'état ;
  \item les points où une preuve est traduite ou réutilisée.
\end{enumerate}

Ajoutez trois colonnes : \emph{qui émet}, \emph{qui vérifie}, \emph{quelle
frontière est franchie}.

\hint{Une machine à états d'authentification est un Logical Engine : le bug
apparaît souvent dans une transition d'erreur ou de repli, pas dans le
chemin nominal.}

\chapter{Exercise 01}
\extitle{CVE-2025-33073 : réflexion NTLM}
\turnindir{ex01}
\exnumber{01}
\exfiles{cve-2025-33073.md}
\exauthorize{un éditeur de texte, NVD, avis Synacktiv}
\makeheaderfiles

CVE-2025-33073 est décrite comme une élévation de privilèges dans le client
SMB Windows, liée à une réflexion d'authentification NTLM. Le cas est
intéressant parce que la preuve est valide, mais elle est utilisée sur une
mauvaise frontière.

\subsection*{Mission}

Dans \texttt{cve-2025-33073.md} :

\begin{enumerate}
  \item décrivez le chemin de réflexion sans fournir de commande
        d'exploitation ;
  \item expliquez pourquoi le serveur accepte une preuve produite pour le
        client ;
  \item identifiez l'hypothèse de liaison au canal qui manque ;
  \item proposez une détection défensive (journal, événement, indicateur) ;
  \item expliquez pourquoi un correctif local sur un seul composant peut
        être insuffisant.
\end{enumerate}

\subsection*{Sources}
\begin{itemize}
  \item NVD : \url{https://nvd.nist.gov/vuln/detail/CVE-2025-33073}
  \item Synacktiv, analyse de CVE-2025-33073.
  \item Documentation Microsoft sur la liaison de canal NTLM.
\end{itemize}

\chapter{Exercise 02}
\extitle{CVE-2025-58726 : Ghost SPNs}
\turnindir{ex02}
\exnumber{02}
\exfiles{cve-2025-58726.md}
\exauthorize{un éditeur de texte, NVD, avis Semperis}
\makeheaderfiles

CVE-2025-58726 est associée à des « Ghost SPNs » et à une confusion de
service dans l'authentification Windows. Le mécanisme exact doit être
vérifié dans les avis publics.

\subsection*{Mission}

Dans \texttt{cve-2025-58726.md} :

\begin{enumerate}
  \item définissez ce qu'est un SPN et comment il est utilisé pour choisir
        un service ;
  \item expliquez le risque d'un SPN fantôme ou mal enregistré ;
  \item reliez le cas à la question de frontière de confiance ;
  \item proposez un contrôle de configuration qui réduit l'exposition ;
  \item listez les affirmations que vous n'avez pas pu vérifier.
\end{enumerate}

\subsection*{Sources}
\begin{itemize}
  \item NVD : \url{https://nvd.nist.gov/vuln/detail/CVE-2025-58726}
  \item Semperis, analyse des Ghost SPNs.
\end{itemize}

\chapter{Exercise 03}
\extitle{2026 : chaîne et correctif incomplet}
\turnindir{ex03}
\exnumber{03}
\exfiles{chain.md}
\exauthorize{un éditeur de texte, NVD, avis éditeurs}
\makeheaderfiles

Étudiez quatre CVE de 2026 :

\begin{itemize}
  \item CVE-2026-24294 : port SMB arbitraire ;
  \item CVE-2026-26128 : homoglyph DNS ;
  \item CVE-2026-26119 : contournement EPA dans Windows Admin Center ;
  \item CVE-2026-50343 : Dark Elevator, chargement d'une DLL plantable
        par \texttt{InstallService}.
\end{itemize}

Dark Elevator combine deux défauts logiques : une carte de plugins
inscriptible par un utilisateur normal, et une classe COM Windows dont la
DLL in-process est enregistrée sous \texttt{\%PROGRAMDATA\%}, donc
plantable. Le service SYSTEM charge la DLL dans son propre processus. Le
suivi Project Zero, CVE-2026-66804, montre que le correctif de
\texttt{InstallService} laisse l'enregistrement COM orphelin.

\subsection*{Mission}

Dans \texttt{chain.md}, remplissez un tableau :

\begin{center}
\begin{tabularx}{\linewidth}{|p{0.2\linewidth}|X|X|}
\hline
\textbf{CVE} & \textbf{Hypothèse du correctif} & \textbf{Frontière contournée} \\ \hline
CVE-2026-24294 & & \\ \hline
CVE-2026-26128 & & \\ \hline
CVE-2026-26119 & & \\ \hline
CVE-2026-50343 & & \\ \hline
\end{tabularx}
\end{center}

Puis répondez :

\begin{enumerate}
  \item Ces quatre cas sont-ils des variantes d'une même racine, ou des
        bugs indépendants qui partagent un thème ?
  \item Quel contrôle défensif couvre plusieurs cas à la fois ?
  \item Pourquoi la validation d'un nom ou d'un port ne suffit-elle pas à
        garantir la frontière ?
  \item Pourquoi le correctif de \texttt{InstallService} ne ferme-t-il pas
        l'enregistrement COM orphelin de CVE-2026-66804 ?
  \item Classez CVE-2026-66804 : \emph{incomplete-fix variant} ou
        \emph{architectural scatter} ? Justifiez avec la définition du
        papier.
  \item Proposez trois contrôles défensifs : ACL de la carte de plugins,
        inventaire des enregistrements COM, surveillance des DLL créées
        sous \texttt{ProgramData}.
\end{enumerate}

\begin{warn}
Le suivi Project Zero montre qu'un correctif peut fermer le chemin
d'activation sans supprimer la ressource dangereuse. Ne reproduisez pas la
chaîne de chargement ; travaillez sur la carte des frontières et les
contrôles de configuration.
\end{warn}

\chapter{Exercise 04}
\extitle{Laboratoire et détection}
\turnindir{ex04}
\exnumber{04}
\exfiles{lab.md}
\exauthorize{un éditeur de texte, documentation AD, outils de journalisation}
\makeheaderfiles

Concevez un laboratoire défensif, pas un terrain d'attaque :

\begin{enumerate}
  \item un domaine AD simulé avec des comptes jetables ;
  \item un serveur SMB isolé ;
  \item un client Windows dans un réseau séparé ;
  \item des journaux centralisés ;
  \item des règles de détection pour les réflexions et les SPN suspects.
\end{enumerate}

Dans \texttt{lab.md}, décrivez :

\begin{itemize}
  \item le schéma réseau et les frontières ;
  \item les événements à collecter ;
  \item trois requêtes de détection ;
  \item les faux positifs attendus ;
  \item les critères d'arrêt du laboratoire.
\end{itemize}

\begin{warn}
Ne testez jamais ces techniques sur un domaine qui ne vous appartient pas.
Une preuve d'autorisation écrite est obligatoire avant toute manipulation.
\end{warn}

\chapter{Rendu et évaluation}

\begin{console}
A-EPI-004-windows-auth-chain/
|-- ex00/fsm.md
|-- ex01/cve-2025-33073.md
|-- ex02/cve-2025-58726.md
|-- ex03/chain.md
+-- ex04/lab.md
\end{console}

\begin{center}
\begin{tabularx}{\linewidth}{|l|X|c|}
\hline
\textbf{Critère} & \textbf{Attendu} & \textbf{Points} \\ \hline
FSM & Acteurs, preuves et frontières explicites & 5 \\ \hline
Réflexion & Mécanisme et liaison de canal & 5 \\ \hline
SPN & Rôle du SPN et contrôle défensif & 4 \\ \hline
Chaîne 2026 & Dark Elevator, correctif incomplet, sources & 4 \\ \hline
Laboratoire & Isolé, détections et critères d'arrêt & 4 \\ \hline
\end{tabularx}
\end{center}

\chapter{Ressources par exercice}

\subsection*{Exercice 00 --- Carte de la machine à états}
\begin{itemize}
  \item Microsoft, NTLM :
        \url{https://learn.microsoft.com/en-us/windows-server/security/kerberos/ntlm-overview}
  \item Microsoft, Kerberos :
        \url{https://learn.microsoft.com/en-us/windows-server/security/kerberos/kerberos-authentication-overview}
  \item Microsoft, SMB :
        \url{https://learn.microsoft.com/en-us/windows-server/storage/file-server/file-server-smb-overview}
  \item Papier, section 16.1.
  \item À chercher : \texttt{NEGOTIATE}, \texttt{CHALLENGE},
        \texttt{AUTHENTICATE}, \texttt{AS-REQ}, \texttt{TGS-REQ}.
\end{itemize}

\subsection*{Exercice 01 --- CVE-2025-33073}
\begin{itemize}
  \item NVD : \url{https://nvd.nist.gov/vuln/detail/CVE-2025-33073}
  \item MSRC :
        \url{https://msrc.microsoft.com/update-guide/vulnerability/CVE-2025-33073}
  \item Synacktiv, analyse NTLM reflection :
        \url{https://www.synacktiv.com/sites/default/files/2026-05/bypassing_authentication_reflections_mitigations_for_system_shells.pdf}
  \item Papier, section 16.2.
  \item À chercher : \texttt{mrxsmb.sys}, \texttt{CMTI}, \texttt{SPN},
        \texttt{channel binding}, \texttt{loopback check}.
\end{itemize}

\subsection*{Exercice 02 --- CVE-2025-58726}
\begin{itemize}
  \item NVD : \url{https://nvd.nist.gov/vuln/detail/CVE-2025-58726}
  \item MSRC :
        \url{https://msrc.microsoft.com/update-guide/vulnerability/CVE-2025-58726}
  \item Semperis, Ghost SPNs :
        \url{https://www.semperis.com/blog/exploiting-ghost-spns-and-kerberos-reflection}
  \item À chercher : \texttt{Ghost SPN}, \texttt{service account},
        \texttt{KDC}, \texttt{ticket}.
\end{itemize}

\subsection*{Exercice 03 --- Chaîne 2026 et correctif incomplet}
\begin{itemize}
  \item NVD : CVE-2026-24294, CVE-2026-26128, CVE-2026-26119,
        CVE-2026-50343, CVE-2026-66804.
  \item MSRC :
        \url{https://msrc.microsoft.com/update-guide/}
  \item Calif, Dark Elevator :
        \url{https://blog.calif.io/p/dark-elevator-windows-install-service}
  \item Project Zero, dangling COM :
        \url{https://projectzero.google/2026/09/windows-dangling-com.html}
  \item Synacktiv, PDF auth reflections (même lien que l'exercice 01).
  \item À chercher : \texttt{InstallService}, \texttt{StaticPluginMap},
        \texttt{CrossDevice}, \texttt{ProgramData}, \texttt{COM orphan}.
\end{itemize}

\subsection*{Exercice 04 --- Laboratoire et détection}
\begin{itemize}
  \item Microsoft, AD DS :
        \url{https://learn.microsoft.com/en-us/windows-server/identity/ad-ds/active-directory-domain-services}
  \item Microsoft, audit events :
        \url{https://learn.microsoft.com/en-us/windows/security/threat-protection/auditing/advanced-security-audit-policy}
  \item Papier, sections 16.2 et 16.3.
  \item À chercher : \texttt{event 4624}, \texttt{event 4768},
        \texttt{event 4769}, \texttt{SPN}, \texttt{loopback}.
\end{itemize}

\chapter{Sources}
\begin{itemize}
  \item Phrack, \emph{Security is failure: Every Patch Is a Guess},
        sections 9 et 10.
  \item NVD : CVE-2025-33073, CVE-2025-58726, CVE-2026-24294,
        CVE-2026-26128, CVE-2026-26119, CVE-2026-50343,
        CVE-2026-66804.
  \item Calif, \emph{Dark Elevator: Windows Install Service Local Privilege
        Escalation} :
        \url{https://blog.calif.io/p/dark-elevator-windows-install-service}
  \item Project Zero, \emph{Dangling COM Object Registrations} :
        \url{https://projectzero.google/2026/09/windows-dangling-com.html}
  \item Synacktiv, \emph{The Gift That Keeps on Giving} :
        \url{https://www.synacktiv.com/sites/default/files/2026-05/bypassing_authentication_reflections_mitigations_for_system_shells.pdf}
  \item Documentation Microsoft sur NTLM, Kerberos, SMB, EPA et COM.
  \item Analyses Semperis.
\end{itemize}

\end{document}
